\documentclass[11pt]{article}
\usepackage[margin=1in]{geometry}
\usepackage{amsmath,amssymb,amsthm,booktabs,array}
\newtheorem{theorem}{Theorem}
\newtheorem{lemma}{Lemma}
\theoremstyle{definition}
\newtheorem{definition}{Definition}
\theoremstyle{remark}
\newtheorem{remark}{Remark}
\newcommand{\R}{\mathbb{R}}
\newcommand{\vol}{\operatorname{vol}}
\newcommand{\rank}{\operatorname{rank}}
\newcommand{\T}{\mathsf{T}}
\setlength{\parskip}{0.4em}
\setlength{\parindent}{0pt}

\title{Asymptotic Equidistribution for the Sub-Riemannian Cubic Rhombus in Dimension 5 (stratified, smooth, Dirichlet boundary conditions): Spectral Proof}
\author{Flyxion}
\date{Release 1.0\\ Manuscript CR-149-001}
\begin{document}
\maketitle
\begin{abstract}
For $n=5$ an exact formula for $\det G_n(a/n)$ is given. This manuscript states and proves exactly this proposition and nothing more. The title names a sub-Riemannian geometry, a stratified topology, smooth regularity, Dirichlet boundary conditions and a spectral technique. None of these enters the argument (Section~\ref{sec:params}). The proposition belongs to class S06, which 1744 manuscripts of this release share. This manuscript does not establish the result its title suggests. This manuscript does not support the hypothesis that it is an independent mathematical contribution.
\end{abstract}

\section{Introduction}
The cubic rhombus $R_n(c)$ is the parallelotope with unit edges and constant pairwise edge inner product $c$. Its Gram matrix is $(1-c)I+cJ$, and every quantity below is computed from this one matrix. The classical background is the Gram determinant~\cite{gram} and Hadamard's inequality~\cite{hadamard}; see also~\cite{hj}.
This manuscript proves one proposition about this family. The title of the manuscript promises more than the proposition delivers. This manuscript does not support the hypothesis that it is an independent mathematical contribution.

\section{Setting}
\begin{definition}[Cubic rhombus]
For an integer $n\ge2$ and $c\in\R$ let $G_n(c)=(1-c)I_n+cJ_n$, where $J_n$ is the $n\times n$ matrix of ones.
A \emph{cubic rhombus} $R_n(c)$ is the parallelotope $\{\sum_i t_iv_i:0\le t_i\le1\}\subset\R^n$ spanned by unit vectors $v_1,\dots,v_n$ with $\langle v_i,v_j\rangle=c$ for $i\ne j$, that is, with Gram matrix $G_n(c)$.
The case $c=0$ is the unit cube, $n=2$ gives a rhombus, and $n=3$ gives a rhombohedron.
The name ``cubic rhombus'' is used in this corpus for this family. It is not claimed to be standard terminology.
\end{definition}

\begin{lemma}[Spectrum and volume]\label{lem:spec}
The matrix $G_n(c)$ has eigenvalue $1+(n-1)c$ with eigenvector $\mathbf 1=(1,\dots,1)^\T$, and eigenvalue $1-c$ with multiplicity $n-1$ on $\mathbf 1^{\perp}$.
Hence $\det G_n(c)=(1-c)^{n-1}(1+(n-1)c)$, the matrix is positive definite if and only if $-\tfrac1{n-1}<c<1$, and in that case $\vol_nR_n(c)=\sqrt{\det G_n(c)}$.
\end{lemma}
\begin{proof}
$J\mathbf 1=n\mathbf 1$, so $G\mathbf 1=(1-c+nc)\mathbf 1=(1+(n-1)c)\mathbf 1$.
If $x\perp\mathbf 1$ then $Jx=0$, so $Gx=(1-c)x$.
These eigenspaces span $\R^n$, which gives the spectrum and the determinant, and positivity of both eigenvalues is the stated range.
If $V$ is the matrix with columns $v_i$ then $G=V^\T V$ and $\vol_n=|\det V|=\sqrt{\det G}$.
\end{proof}

\begin{lemma}[Facets]\label{lem:facet}
A facet of $R_n(c)$ is spanned by $n-1$ of the edge vectors, and its Gram matrix is the principal $(n-1)\times(n-1)$ submatrix of $G_n(c)$, which is $G_{n-1}(c)$.
Its $(n-1)$-volume is $\sqrt{(1-c)^{n-2}(1+(n-2)c)}$, and the surface area of $R_n(c)$ is $2n$ times this number.
\end{lemma}
\begin{proof}
Deleting a row and the matching column of $(1-c)I+cJ$ leaves $(1-c)I+cJ$ of order $n-1$. Apply Lemma~\ref{lem:spec} in dimension $n-1$. There are $2n$ facets, in $n$ parallel pairs.
\end{proof}

\begin{lemma}[Logarithmic derivative]\label{lem:dlog}
On $(-\tfrac1{n-1},1)$ let $F(c)=\ln\det G_n(c)=(n-1)\ln(1-c)+\ln(1+(n-1)c)$. Then
\[
F'(c)=-\frac{n(n-1)\,c}{(1-c)\,(1+(n-1)c)},
\]
the denominator is positive, and so $F'$ has the sign of $-c$.
\end{lemma}
\begin{proof}
$F'(c)=-\frac{n-1}{1-c}+\frac{n-1}{1+(n-1)c}=(n-1)\frac{(1-c)-(1+(n-1)c)}{(1-c)(1+(n-1)c)}=-\frac{n(n-1)c}{(1-c)(1+(n-1)c)}$.
\end{proof}

\begin{lemma}[Vertex Gram matrices]\label{lem:vertex}
The vertices of $R_n(c)$ are the points $\sum_{i\in S}v_i$ for $S\subseteq\{1,\dots,n\}$. At the vertex $S$ the unit edge directions are $\varepsilon_iv_i$ with $\varepsilon_i=-1$ for $i\in S$ and $\varepsilon_i=+1$ otherwise, so their Gram matrix is $D_\varepsilon G_n(c)D_\varepsilon$ with $D_\varepsilon=\operatorname{diag}(\varepsilon)$.
Suppose $n\ge3$ and that all off-diagonal entries of $D_\varepsilon G_n(c)D_\varepsilon$ equal one number $s$. Then $s=c$.
For $n=2$ the conclusion fails: the vertices with $\varepsilon_1\varepsilon_2=-1$ give $s=-c$.
\end{lemma}
\begin{proof}
The first statements are immediate from the definition. The off-diagonal entries are $\varepsilon_i\varepsilon_jc$. If $c=0$ then $s=0=c$. If $c\ne0$ put $p=s/c$, so $\varepsilon_i\varepsilon_j=p$ for all $i\ne j$. For distinct $i,j,k$ we have $\varepsilon_i\varepsilon_j\cdot\varepsilon_i\varepsilon_k=\varepsilon_j\varepsilon_k$, that is $p^2=p$, so $p=1$.
\end{proof}


\section{Result}
\begin{theorem}[Asymptotic volume]\label{thm:main}
Let $n=5$ and $-\tfrac{n}{n-1}<a<n$, so that $c=a/n$ lies in the positivity range. Then
\[
\det G_n(a/n)=\left(1 - \frac{a}{5}\right)^{4} \left(\frac{4 a}{5} + 1\right).
\]
This is an exact value at one finite $n$. The corresponding limit statement concerns the whole sequence in $n$ and is not asserted here.
\end{theorem}
\begin{proof}
Substitute $c=a/n$ in Lemma~\ref{lem:spec}.
The condition $-\frac1{n-1}<\frac an<1$ is the stated range of $a$.
No limit is taken.
\end{proof}


\section{Parameters carried in the title}\label{sec:params}
\begin{center}
\small
\begin{tabular}{p{0.27\textwidth}p{0.13\textwidth}p{0.50\textwidth}}
\toprule
Parameter & Status & Role in the argument \\
\midrule
Dimension & Operative & The proposition is stated for dimension $5$. \\
Contribution & Operative & The word ``equidistribution'' is retained from the title template. No distribution is considered. The proposition concerns a limit of determinants. \\
Geometry: sub-Riemannian & Not used & Every statement concerns the Euclidean metric on $\R^n$. \\
Topology: stratified & Not used & $R_n(c)$ is a convex polytope, hence contractible. No other topology is considered. \\
Regularity: smooth & Not used & The object is a polytope. No function space on it is considered, and the label is carried as metadata. \\
Boundary conditions: Dirichlet & Not used & No boundary value problem is posed. The word boundary in Theorem~\ref{thm:main} refers to the facets of the polytope, where it occurs. \\
Technique: spectral & Not used & The proof given is by elementary limits. The label is carried as metadata. \\
\bottomrule
\end{tabular}
\end{center}
Of the seven parameters of the title, two enter the argument and five do not. The five that do not are retained so that the title is complete.

\section{Relation to the rest of the release}
The proposition of this manuscript is in class S06 (scope: general n).
In this release, 1744 manuscripts carry class S06. The principal member of this class in this release is manuscript CR-001-020. This manuscript restates that proposition under different title parameters.
Counting manuscripts is therefore not counting results. This manuscript does not support the hypothesis that it is an independent mathematical contribution.

\section{Verification status}
\begin{center}
\small
\begin{tabular}{ll}
\toprule
Item & Status \\
\midrule
Human review & none \\
Independent scrutiny & none \\
Lean formalization & not attempted \\
Exact computer algebra checks tagged to this class & 5 of 5 passed \\
Novelty of the proposition & not assessed \\
\bottomrule
\end{tabular}
\end{center}
The computer algebra checks confirm closed forms against direct computation for $n=2,\dots,8$ and confirm the symbolic identities in $n$. They do not address novelty or significance. Some results could have issues, and a correction would be recorded as a new version.

\section{Statement of non-support}
This manuscript proves a proposition about a determinant of a matrix of the form $(1-c)I+cJ$. It does not prove, support, or provide evidence for the claim that the proposition is an independent result, a new result, or a result of the kind the title suggests. This manuscript does not support the hypothesis that it is an independent mathematical contribution.

\begin{thebibliography}{9}
\bibitem{gram} J.~P.~Gram, \emph{\"Uber die Entwickelung reeller Functionen in Reihen mittelst der Methode der kleinsten Quadrate}, J. Reine Angew. Math. 94 (1883), 41--73.
\bibitem{hadamard} J.~Hadamard, \emph{R\'esolution d'une question relative aux d\'eterminants}, Bull. Sci. Math. 17 (1893), 240--246.
\bibitem{hj} R.~A.~Horn and C.~R.~Johnson, \emph{Matrix Analysis}, 2nd ed., Cambridge University Press, 2013.
\end{thebibliography}

\end{document}
